\scanpage{46}
% L1
où chaque
\[
\deleted{\unsure{\beta_{j}}}\gamma_{j}(X)\in\mathbb Z^{\sigma_j(X)}
\qquad 0\leq j\leq n
\]
Ils sont obtenus en prenant les $\beta_k(X)$ de Serre [composantes de poids purs de $\mathrm{cl}(X)$] pour $k\geq n$, et \deleted{les composant}\inserted{en notant} que un tel $\beta_k(X)$ se met sous la forme $\xi^{k-n}\beta'_{\deleted{\unsure{k}}2n-k}(X)$, où $\beta'_{\deleted{\unsure{k}}2n-\unsure{n|k}}(X)\in M^+_{2n-k}(K)$ ; alors $\gamma_{2n-\unsure{n|k}}(X)$ s'obtient en prenant la composante \uline{bipure} de $\beta'_{2n-k}(X)$ \dots\ .
% slanted
\margin{aux signes près !}

% L10
Par construction $=$, les invariants ainsi obtenus sont des invariants $n$-birationnels de $X$. Si $X$ est projective \deleted{non singulière}\inserted{lisse sur $K$}, alors $\beta'_j(X)$ pour $0\leq j\leq n$ n'est autre (par \uuline{Lefschetz} \dots) que la classe des $H^j(X,\mathbb Q(0))$. Donc notre construction consiste à prendre la composante $j$-bipure des $H^j(X,\mathbb Q(0))$. Mais cela n'est autre aussi
% archive sheet number: 45
